%!TEX root = ../main.tex

\chapter{Fase 4: integrazione del ruolo della tutor}

\section{La supervisione del tirocinio}

La Prof.ssa Alessandra Lumini ha svolto il ruolo di tutor del tirocinio, accompagnando il percorso di sviluppo del progetto \builder{} e fornendo un riferimento per l'impostazione generale del lavoro. Il suo ruolo deve essere presentato nella relazione finale come una supervisione metodologica e progettuale, non come una semplice presenza formale. La tutor ha contribuito a orientare il lavoro verso obiettivi coerenti con il contesto universitario, aiutando a mantenere attenzione non solo alla realizzazione tecnica, ma anche alla qualita del risultato.

Nel corso del tirocinio, la supervisione ha permesso di considerare il progetto come un'attivita formativa completa. \builder{} non e stato trattato soltanto come un'applicazione da implementare, ma come uno strumento da rendere comprensibile, verificabile e utile per la progettazione di basi di dati. Questo ha richiesto di collegare le funzionalita tecniche alla chiarezza concettuale e alla leggibilita degli schemi.

Il contributo della tutor puo essere valorizzato anche in relazione alla definizione delle priorita. In un progetto con molte componenti, come editor grafico, traduzione logica, export, salvataggio, localizzazione e test, e importante capire quali aspetti incidono maggiormente sulla qualita complessiva. La supervisione della Prof.ssa Lumini ha aiutato a mantenere un orientamento ordinato, evitando che lo sviluppo si riducesse a una somma di funzionalita scollegate.

\section{Indicazioni metodologiche e progettuali}

Le indicazioni metodologiche della tutor hanno riguardato soprattutto la necessita di mantenere il progetto chiaro e concettualmente corretto. Nel caso di un'applicazione per schemi ER e logici, la correttezza non e un aspetto secondario: ogni scelta grafica o tecnica deve rispettare il significato degli elementi rappresentati. Entita, relazioni, cardinalita, identificatori e generalizzazioni devono essere gestiti in modo coerente, perche da essi dipende la successiva trasformazione verso lo schema logico.

Un secondo tema riguarda la leggibilita grafica. Gli schemi di basi di dati possono diventare rapidamente complessi, soprattutto quando includono molte entita, attributi, relazioni e vincoli. Per questo motivo, i suggerimenti della tutor possono essere collegati alla cura del layout, alla disposizione degli elementi, alla leggibilita delle label e alla riduzione delle sovrapposizioni. La chiarezza visiva ha valore didattico, perche consente allo studente di interpretare correttamente lo schema e di discutere le scelte progettuali.

Un ulteriore aspetto riguarda la semplicita d'uso. Un'applicazione didattica deve permettere all'utente di concentrarsi sui concetti di modellazione, senza essere ostacolato da un'interfaccia confusa. La revisione progressiva di toolbar, pannelli, modali, messaggi e flussi utente puo quindi essere letta anche come risposta alla necessita di rendere lo strumento piu accessibile e ordinato.

Infine, la tutor ha rappresentato un riferimento per mantenere attenzione alla qualita del codice e alla manutenibilita. Un progetto destinato a evolvere deve avere una struttura chiara, test adeguati e una separazione ragionevole tra interfaccia, logica di dominio e tipi condivisi. Questo approccio ha favorito una crescita piu controllata dell'applicazione.

\section{Impatto dei suggerimenti sul progetto}

\noindent\begin{tabularx}{\textwidth}{>{\bfseries}p{4cm}X}
Indicazione della tutor & Impatto sul progetto \\
\hline
Attenzione alla leggibilita degli schemi & Miglioramento di layout, label, disposizione grafica, routing e gestione delle sovrapposizioni. \\
Chiarezza dell'interfaccia & Maggiore cura di toolbar, pannelli, modali, command menu, loading screen e flussi utente. \\
Correttezza concettuale & Controllo piu attento di entita, relazioni, cardinalita, identificatori, generalizzazioni e traduzione logica. \\
Semplicita d'uso & Riduzione del rumore visivo e maggiore attenzione a messaggi, scorciatoie, onboarding e responsive. \\
Manutenibilita & Organizzazione modulare del codice, separazione tra UI e logica di dominio, uso di tipi condivisi e test. \\
Valore didattico & Descrizione piu chiara dei flussi ER/logico e attenzione al percorso di apprendimento dell'utente. \\
Revisione progressiva & Miglioramento incrementale del progetto, con correzione di bug, consolidamento delle funzionalita e aggiornamento della documentazione. \\
\end{tabularx}

\section{Competenze sviluppate grazie al confronto con la tutor}

\subsection{Competenze tecniche}

Il confronto con la tutor ha favorito lo sviluppo di competenze tecniche legate alla progettazione di uno strumento coerente con il dominio delle basi di dati. L'attenzione alla correttezza degli schemi ha richiesto di comprendere meglio il rapporto tra modello ER, modello logico, chiavi, vincoli e trasformazioni. Inoltre, la necessita di rendere leggibili gli schemi ha portato a lavorare su layout, canvas SVG, export e interfaccia utente.

\subsection{Competenze metodologiche}

Dal punto di vista metodologico, la supervisione ha aiutato a procedere per priorita, verificando progressivamente le scelte adottate. Il tirocinio ha richiesto di distinguere tra problemi urgenti, miglioramenti strutturali e possibili sviluppi futuri. Questa capacita di pianificazione e revisione e particolarmente importante in un progetto software, dove ogni modifica puo avere effetti su piu parti del sistema.

\subsection{Competenze comunicative}

Il confronto con la tutor ha contribuito anche allo sviluppo di competenze comunicative. Spiegare le scelte progettuali, motivare le soluzioni adottate e descrivere i limiti del lavoro sono attivita necessarie per trasformare un progetto tecnico in una relazione universitaria. La capacita di comunicare in modo chiaro il rapporto tra funzionalita, obiettivi e risultati rappresenta quindi una parte significativa dell'esperienza formativa.

\section{Paragrafo pronto da inserire nella relazione finale}

Durante il tirocinio, la Prof.ssa Alessandra Lumini ha svolto il ruolo di tutor, supervisionando il percorso di sviluppo del progetto \builder{} e fornendo indicazioni utili per orientare le scelte progettuali. Il suo contributo ha aiutato a mantenere il lavoro focalizzato non soltanto sulla realizzazione tecnica delle funzionalita, ma anche sulla chiarezza concettuale, sulla qualita dell'interfaccia e sull'utilita didattica dello strumento. In un'applicazione dedicata alla progettazione di basi di dati, infatti, la correttezza degli schemi ER e logici rappresenta un requisito fondamentale: entita, relazioni, attributi, cardinalita, identificatori e generalizzazioni devono essere gestiti in modo coerente per supportare una trasformazione logica comprensibile e verificabile. Le indicazioni della tutor hanno contribuito a dare importanza anche alla leggibilita grafica degli schemi, alla disposizione degli elementi sul canvas e alla semplicita dei flussi utente, aspetti essenziali per rendere l'applicazione adatta a un contesto universitario. Il confronto con la Prof.ssa Lumini ha inoltre favorito una maggiore attenzione alla manutenibilita del codice, alla separazione tra interfaccia e logica di dominio e alla necessita di verificare le parti piu delicate tramite test. Nel complesso, la supervisione ricevuta ha rappresentato un elemento importante del percorso formativo: ha permesso di affrontare il progetto con maggiore metodo, di definire meglio le priorita e di sviluppare competenze tecniche, metodologiche e comunicative. Grazie a questo confronto, il tirocinio e stato vissuto non solo come attivita di sviluppo software, ma come esperienza di progettazione, revisione e crescita professionale.
